% =====================================================================
% c00.tex —— 第 1 节「はじめに」/ 第 2 节「ミラー対称性」（p001–p005）
% 底本：植田一石《ホモロジー的ミラー対称性》讲义录
% 公式一律以 work/pages/pNNN.png 页面图像为准逐符号核对。
%
% [Tutor Note: BRIEF 记 c00 有 10 个编号公式，该数来自文字层命中（含正文对同一编号的交叉引用）。按页面图像核对，p001–p005 上实际编号公式为 (2.1)–(2.6) 共 6 条；此处按实际条数标注，未凑数、未给无编号公式加 tag。]
% [Tutor Note: 原文以「予想 N.M.」行内粗体形式给出猜想陈述（并非独立环境）；preamble.tex 未提供 conjecture / 予想 环境，故按 STYLE_GUIDE §2 的「行内粗体小标题」体例用 \para{猜想 N.M.} 保留原文编号。]
% [Tutor Note: 原文 p001 第 2 段作「Lagarnge 多様体」，当为「Lagrange」之误；译文按 Lagrange 书写，未改原意。]
% =====================================================================

\section{引言}

镜像对称是源自超弦理论的一种数学现象，指某个空间的辛几何与另一个空间的复几何之间存在不可思议的关系。它最初只以 Calabi–Yau 流形为对象，但其适用范围逐渐扩大，如今甚至有一种势头，让人觉得对所有的复流形与辛流形，在某种意义下镜像都存在。

以对各种与镜像对称相关的现象给出概念性理解为目标、由 Kontsevich 于 1994 年提出的猜想，就是同调镜像对称。它断言：在构成镜像对（mirror pair）的一对流形中，相干层的导出范畴与 Lagrange 子流形所构成的 Fukaya 范畴的导出范畴是等价的。这一猜想提供了把 A∞ 范畴看作空间并研究其几何学这一全新的视角，蕴含着从根本上颠覆我们空间概念的可能性。在提出 20 多年之后，同调镜像对称不仅把复几何与辛几何联系起来，还与完全可积系统、纽结理论、几何 Langlands 对应、簇代数等广泛的数学相关联，成长为一个巨大的领域。这里，在简要回顾镜像对称的历史之后，我想对同调镜像对称的几个侧面作一概观。\footnote{本稿是在东京大学大学院数理科学研究科 2015 年度冬学期讲义录的基础上增补、修订而成的。}

\section{镜像对称}

镜像对称的源流可以上溯到弦理论\footnote{弦理论的教科书例如有 [GSW87, Pol98b]。}的环面紧致化\footnote{弦理论中的环面紧致化，是指取 $T$ 为 $k$ 维实环面、$M$ 为 $n-k$ 维流形时，在直积 $T\times M$ 上考虑弦理论；它与 [AMRT10] 所处理的数学中的环面紧致化没有直接关系。这里 $n$ 是自然数，通常取弦理论的临界维数（10 或 26）。}的研究过程中所发现的一个认识：弦以与点粒子完全不同的方式观看这个世界 [KY84, SS86]。对在半径为 $2\pi R$ 的圆周 $T\cong\R/2\pi R\Z$ 上运动的自由粒子，与时间无关的 Schrödinger 方程为
\begin{equation}
    -\frac{\partial^2}{\partial x^2}\Psi(x) = E \Psi(x) \tag{2.1}
\end{equation}
其解用动量 $m\in\Z$ 写成 $\Psi(x)=\exp\bigl(\sqrt{-1}\,mx/R\bigr)$，此时的能量由 $E=m^2/R^2$ 给出。另一方面，为了描述 $T$ 上经典弦的运动，考虑以泛函
\begin{equation}
    \operatorname{Map}(S^1,T)\to\R,\quad
    \gamma\mapsto\int_{S^1}\left(\frac{d\gamma}{d\theta}\right)^2\frac{d\theta}{2\pi} \tag{2.2}
\end{equation}
作为能量，则对给定的卷绕数 $w\in\Z\cong\pi_1(T)$，使能量最小的位形为 $\gamma(\theta)=Rw\theta$，此时的能量由 $R^2w^2$ 给出。在 $T$ 上运动的量子弦同时具有动量与卷绕数，其能量由
\begin{equation}
    E=\frac{m^2}{R^2}+R^2w^2 \tag{2.3}
\end{equation}
给出。这里，若在把半径 $R$ 换成 $1/R$ 的同时交换动量与卷绕数，则可以看出弦的能量保持不变。这一论证由于本质上非相对论、以及描述弦的状态时除动量与卷绕数之外还需要无限多个指定弦振动模式的自然数等原因而不准确，但经更精细的考察可知，半径为 $R$ 的圆周上的弦理论与半径为 $1/R$ 的圆周上的弦理论是不可区分的。这被称为 T 对偶性，其中以下两点很重要：

\begin{itemize}
    \item 这是点粒子的理论中所没有的、弦理论特有的现象。
    \item 这是经典理论中所没有的量子现象。
\end{itemize}

进而，在 T 对偶变换 $R\mapsto1/R$ 的不动点 $R=1$ 处，会发生称为规范对称性增强（enhancement）的现象，这对于用无限维 Lie 环的顶点算子代数来实现而言是本质的\footnote{关于弦理论的环面紧致化与 T 对偶性的入门性记述，参见例如 [Pol98a, Chapter 8]。}。

更一般地，取 $n$ 维 Euclid 空间 $\R^n$，用与 $\Z^n$ 同构的离散子群 $N\subset\R^n$ 作商所得的环面记为 $T:=\R^n/N$，则在 $T$ 上运动的弦，除了取值为 $T$ 的特征标群 $M:=\Hom(T,S^1)\cong\Hom_{\Z}(N,\Z)$ 的动量之外，还具有取值为 $N\cong\pi_1(T)$ 的卷绕数；弦的动力学在同时交换动量与卷绕数、并交换环面 $T=\R^n/N$ 与对偶环面 $\check{T}=(\R^n)^{\vee}/M$ 的操作下不变。换言之，对弦理论而言，环面 $T$ 与对偶环面 $\check{T}$ 是无法分辨的。

% [Tutor Note: 原文「混成弦理論」为 heterotic string theory，大陆通行译名为「杂化弦」，此处按通行译名处理（首次出现附原文 heterotic）；glossary.tsv 未收该词，若需与其余 chunk 统一请 Lead 裁定。]

超弦理论只存在于 10 维，分为 I 型超弦理论、IIA 型超弦理论、IIB 型超弦理论、$SO(32)$ 杂化弦理论（heterotic）、$E_8\times E_8$ 杂化弦理论 5 类。其中之一的 $E_8\times E_8$ 杂化弦理论，若在实 4 维时空 $\R^4$ 与复 3 维 Calabi–Yau 流形 $M$ 的直积 $\R^4\times M$ 上考虑，则已知它给出接近粒子物理标准模型的模型，于是 3 维 Calabi–Yau 流形上的弦理论研究就盛行起来\footnote{这里 3 维 Calabi–Yau 流形是指 Ricci 曲率为零的 3 维 Kähler 流形。这与说全纯群包含于 $SU(3)$ 中的实 6 维 Riemann 流形是同一回事。注意，根据 Yau 对 Calabi 猜想的解决，具有平凡典范丛的紧 Kähler 流形在相差常数倍的意义下具有唯一的 Ricci 平坦度量。全纯群通过协变常数的旋量与时空的超对称性的大小相关。为了使 $\R^4\times M$ 上的 $E_8\times E_8$ 杂化弦理论给出接近粒子物理标准模型的模型，全纯群必须恰好与 $SU(3)$ 一致，但本论说不作这一假设。}。

超弦理论的数学上严格的构造尚不可知，但人们期待：在 3 维 Calabi–Yau 流形上考虑杂化弦理论会给出 2 维的 $(0,2)$ 超共形场论，而在 3 维 Calabi–Yau 流形上考虑 II 型超弦理论会给出 2 维的 $(2,2)$ 超共形场论。

2 维的 $(2,2)$ 超共形场论中有一种操作，把与超共形代数所含的左行（left-moving；数学文献中多把左行与右行分别称为全纯与反全纯）$U(1)$ 流相伴的标量场的符号反转。于是，对任意 Calabi–Yau 流形 $M$，由 $M$ 上的 II 型超弦理论所确定的超共形场论经此操作后得到的超共形场论，是否再次来自某个适当的 Calabi–Yau 流形 $W$，就成了自然的问题。当这一问题的答案为肯定时，称 $W$ 为 $M$ 的镜像流形（mirror manifold）。

由于超弦理论的数学上严格的定义尚不可知，要把镜像对称数学地定式化，就需要考虑比超共形场论更粗的对象，即对任意 Calabi–Yau 流形都数学地严格确定的某种东西。这种较粗信息的典型例子是 Hodge 数，用它可以把称为位相镜像对称（topological mirror symmetry）的猜想表述如下：

\para{猜想 2.1.} 对任意 3 维 Calabi–Yau 流形 $M$，存在某个 3 维 Calabi–Yau 流形 $W$，使得对任意 $0\le p,q\le3$，$M$ 与 $W$ 的 Hodge 数之间
\begin{equation}
    h^{p,q}(M)=h^{p,3-q}(W) \tag{2.4}
\end{equation}
这样的关系成立。

特别地有 $h^{1,1}(M)=h^{1,2}(W)$，但须注意左边是 $M$ 的 Kähler 类的形变自由度的维数，右边是 $W$ 的复结构的形变自由度的维数。任意 Calabi–Yau 流形都具有把 Kähler 类乘以常数的自由度，因此 $h^{1,1}(M)$ 总是 1 以上；但存在没有形变的 Calabi–Yau 流形，所以 $h^{1,2}(W)$ 可以等于 0。因此，猜想 2.1 照原样并不成立。另外，若 Calabi–Yau 流形是单连通的，则其 Euler 数由 $\chi=2(h^{1,1}-h^{1,2})$ 给出，因此镜像对称把 Euler 数的符号反转。

作为 Calabi–Yau 流形的例子，立刻能想到的是加权射影空间中超曲面的奇点消解。按朴素的想法，这样得到的 Calabi–Yau 流形似乎具有来自齐次多项式个数的大 $h^{1,2}$ 和来自加权射影空间 Picard 群的小 $h^{1,1}$。所以当弦理论家把作为 4 维加权射影空间中超曲面得到的 6000 个 3 维 Calabi–Yau 流形的 Hodge 数绘出图来时，其中发现了许多位相镜像对（即满足 (2.4) 的 Calabi–Yau 流形对 $(M,W)$），这一事实受到代数几何学家的惊喜欢迎\footnote{参见 [CLS90, Fig. 1] 中令人印象深刻的图。}。另外，他们计算 Hodge 数所用的 Calabi–Yau/Landau–Ginzburg 对应与 Vafa 的公式在当时尚未得到数学上的正当化，而为理解它们所作的努力成了此后发展的原动力\footnote{关于 Calabi–Yau/Landau–Ginzburg 对应的数学研究，例如参见 [Orl09, CR10, CR11, CIR14] 及其参考文献。特别地，[Orl09] 与同调镜像对称关系很深。另外，Vafa 的公式产生了 3 维 McKay 对应 [Rei, Rei02]、弦理论式 Hodge 数 [Bat98]、轨形量子上同调 [CR02, AGV08]、Motif 积分 [DL99, Cra04] 等理论。}。关于位相镜像对称最著名的结果是 Batyrev 等人 [Bat94, BD96, BB96] 的工作，他们证明了：对于作为 Gorenstein 环面 Fano 流形上 nef 线丛的完全交所得的 Calabi–Yau 流形，猜想 2.1 成立。

% [Tutor Note: 原文「地誌学」对应英文 geography（曲面地理学），此处按字面译作「地志学」；国内文献或作「地理学」，请审校统一。]

位相镜像对称是一般型曲面的地志学 [Per87] 在 Calabi–Yau 流形中的类似物，至今仍是持续被研究的深刻问题，但在与镜像对称相关的诸猜想中是最弱的。使镜像对称引起数学家极大关注的是 [CdlOGP91] 所给出的惊人预言：$\PP^4$ 中 5 次超曲面上的有理曲线的计数与其镜像流形的周期有关。

计算给定代数流形上给定次数的有理曲线条数，属于称为计数几何（enumerative geometry）的领域，是一个经典问题。$\PP^3$ 中光滑 3 次曲面有 27 条直线这一事实很有名，而关于 $\PP^4$ 中的 5 次超曲面，有 Clemens [Cle87] 的如下猜想：

\para{猜想 2.2.} 对任意正整数 $d$，$\PP^4$ 的一般 5 次超曲面上的 $d$ 次光滑有理曲线至多只有有限条存在。

猜想 2.2 受到广泛关注，但目前成否不明。另外，已知它与 Lang 的猜想不相容 [Voi03]。假设猜想 2.2 成立，记 $d$ 次光滑有理曲线的条数为 $n_d$。$n_1=2875$ 是经典的，$n_2=609250$ 已由 [Kat86] 算出，但在 1991 年时 $n_3$ 以后的数尚属未知。[CdlOGP91] 猜想了用镜像流形的周期表示 $n_d$ 的生成函数的公式，并由此对 $d\le10$ 具体计算了导出的 $n_d$。在这些预言中，$n_3=317206375$ 由 [ES95] 证明，$n_4=242467530000$ 由 [Kon95a] 证明。

另一方面，不是对曲线作朴素的计数，而是在稳定映射（stable map）的模空间上的虚拟基本类（virtual fundamental class）中作积分，Gromov–Witten 不变量（Gromov–Witten invariant）的概念经过许多人的努力被表述了出来\footnote{例如参见 [CK99] 第 7 章及其参考文献。}。由此回避了与未解决的猜想 2.2 相伴的困难，从而可以把古典镜像对称（classical mirror symmetry）表述如下：

\para{猜想 2.3.} 对任意 Calabi–Yau 流形 $X$，存在某个 Calabi–Yau 流形的退化族 $\cY\to B$，使得 $X$ 的 Gromov–Witten 不变量的生成函数由 $\cY$ 的周期具体描述。

更准确地说，古典镜像对称被表述为：由 $X$ 的 Gromov–Witten 不变量所确定的 $H^2(X;\C)$ 上的 Hodge 结构的形变，与 $\cY$ 所确定的 $B$ 上的 Hodge 结构的形变之间的同构。也有时表述为：由 $X$ 的 Gromov–Witten 不变量所确定的 Frobenius 流形，与由 $\cY$ 的 Hodge 结构的形变所确定的 Frobenius 流形之间的同构。这一猜想对于作为环面 Fano 流形中 nef 线丛的完全交给出的 Calabi–Yau 流形，已由 Givental [Giv96, Giv98] 与 Lian–Liu–Yau [LLY97] 证明；但其证明基于关于环面作用的同变上同调的局部化，这种关系之所以存在的根本理由仍作为谜团留了下来。

以对镜像对称给出概念性理解为目标、由 Kontsevich [Kon95b] 在 1994 年国际数学家大会上提出的，就是如下的同调镜像对称（homological mirror symmetry）：

\para{猜想 2.4.} 对任意 Calabi–Yau 流形 $X$，存在某个 Calabi–Yau 流形 $Y$，使得在 $X$ 的 Fukaya 范畴的导出范畴与 $Y$ 的相干层的导出范畴之间存在增广三角范畴（enhanced triangulated category）的等价；
\begin{equation}
    D^b\mathfrak{Fuk}\,X\cong D^b\operatorname{coh}Y. \tag{2.5}
\end{equation}

(2.5) 左边出现的范畴的对象是 $X$ 的 Lagrange 子流形，态射是 Floer 上同调；而右边出现的范畴的对象是 $M$ 上的相干层，态射是层的上同调。

% [Tutor Note: 原文此处作「$M$ 上の連接層」，按上下文（右边为 $Y$ 的相干层的导出范畴）应为「$Y$ 上」，疑为笔误；译文照原文保留，请审校确认。]

猜想 2.4 主张在乍看之下毫无关系的这两个范畴之间存在深刻的联系，这一点是惊人的，它被认为是与镜像对称相关的诸猜想中最强的一个。在最简单的椭圆曲线情形，(2.5) 与 $\vartheta$ 函数的双线性关系式密切相关，这一点由 [Kon95b] 指出；在此基础上，[PZ98] 证明了比增广三角范畴的等价更弱的分次范畴（graded category）的等价。仅凭分次范畴的结构无法决定的高次 Massey 积的比较，在 [Pol00] 中讨论。另外，对于更一般的 Abel 流形，[KS01] 与 [Fuk02c] 中得到了部分结果。另一方面，通过与这些完全不同的划时代方法，[Sei15] 证明了 $\PP^3$ 中 K3 超曲面情形的猜想 2.4，并由 [She15, She15] 推广到一般维数射影空间中的 Calabi–Yau 超曲面。这些结果作为特例包含了椭圆曲线情形猜想 2.4 的完全解决。此外，对于作为两条椭圆曲线直积给出的 Abel 曲面，通过归约到各椭圆曲线的情形，猜想 2.4 已被证明 [AS10b]\footnote{这一证明本质地使用了称为 Floer 同调（quilted Floer homology）的理论；关于这一理论参见 [WW10] 等参考文献。}。

Seidel 的手法可以理解为：把关于超曲面的同调镜像猜想归约到关于外围射影空间的同调镜像猜想。所谓外围空间，不只是射影空间，也可以考虑更一般的环面 Fano 流形。与 Calabi–Yau 流形的情形不同，环面 Fano 流形的镜像不是流形而是洛朗多项式；此时同调镜像对称的主张如下：

\para{猜想 2.5（[Kon98, Sei01a]）。} 对任意环面 Fano 流形 $M$，存在某个洛朗多项式 $W$，使得 $M$ 的相干层的导出范畴，与 $W$ 所确定的 Lefschetz 纤维丛的 Fukaya–Seidel 范畴的导出范畴，作为增广三角范畴等价；
\begin{equation}
    D^b\operatorname{coh}M\cong D^b\mathcal{F}(W). \tag{2.6}
\end{equation}

猜想 2.5 对环面 Fano 叠、以及对未必是 Fano 的环面叠也可推广。Lefschetz 纤维丛是 Morse 函数在辛几何中的类似物。对于 Lefschetz 纤维丛，以消没圈（vanishing cycle）为基的自由 Abel 群上以相交数引入二次形式，所得之物称为 Milnor 格，是奇点论中重要的研究对象。对给定的 Lefschetz 纤维丛具体描述其 Milnor 格是困难的问题。Seidel 基于 Kontsevich [Kon98] 的想法，对 Lefschetz 纤维丛定义了这样的 A∞ 范畴：对象是消没圈，态射由 Lagrange 相交 Floer 复形给出 [Sei01b]。这就是 Fukaya–Seidel 范畴（Fukaya–Seidel category）；在其 Grothendieck 群与 Milnor 格同构这一意义上，它给出了 Milnor 格的范畴化（categorification）。猜想 2.5 是关于 Fukaya–Seidel 范畴的极其非平凡的主张：对 $\PP^2$ 与 $\PP^1\times\PP^1$ 已在 [Sei01a] 中、对环面 del Pezzo 曲面已在 [Ued06] 中、对加权射影平面及一般的 del Pezzo 曲面已在 [AKO08, AKO06] 中得到证明。另外，[UY13] 通过使用称为二聚体模型（dimer model）的组合论对象，对把环面 del Pezzo 曲面用环面的任意有限子群作商所得的环面 Fano 叠证明了猜想 2.5。利用这一二聚体模型的理论，还可证明对任意 2 维环面 Fano 叠 $X$，存在某个 Lefschetz 纤维丛 $W$，使得导出范畴的等价 (2.6) 成立 [FU10]\footnote{由二聚体模型构成的 Lefschetz 纤维丛能否用洛朗多项式表示，目前尚不清楚，因此 [FU10] 的主张比猜想 2.5 更弱。}。二聚体模型是 2 维特有的对象，但人们期待存在其高维类似物 [FU14, FLS]。另外，[NU12] 组合 [Sei15] 的手法与 [She11] 的结果证明了 $\PP^4$ 中 5 次超曲面的同调镜像对称，并在过程中得到了对一般维数射影空间的猜想 2.5 的证明。另一方面，虽不是猜想 2.5 本身，但与其密切相关的结果已在 [Abo06, Abo09] 及 [FLTZ11a, FLTZ11b] 中得到。

在 Kontsevich 提出同调镜像对称的翌年，Polchinski 证明了弦理论中所含的称为 D 膜（D-brane）的对象是具有 Ramond–Ramond 荷的 BPS 态 [Pol95]。这是代表超弦理论第二次革命（the second superstring revolution）的成果之一，加速了被称为膜（brane）的扩展对象（extended object）的研究。在这一潮流中，T 对偶性与镜像对称的关系被重新审视，Strominger–Yau–Zaslow [SYZ96] 提出了如下猜想：

\para{猜想 2.6.}
\begin{enumerate}
    \item 任意 Calabi–Yau 流形都具有特殊拉格朗日环面纤维化。
    \item 在这一特殊拉格朗日环面纤维化中，把每条纤维置换成对偶环面，就得到镜像流形。
    \item 同调镜像对称由沿这一特殊拉格朗日环面纤维化的纤维所作的某种 Fourier 变换来实现。
\end{enumerate}

猜想 2.6 一般被称为 SYZ 猜想，其准确定式化与证明是 Calabi–Yau 流形几何中的中心问题之一。特殊拉格朗日子流形是指全纯体积形式的虚部之限制消失的拉格朗日子流形。特殊拉格朗日子流形是校准子流形（calibrated submanifold），从而是 Riemann 几何意义下的极小子流形（minimal submanifold）。由此可知，「是特殊拉格朗日子流形」这一条件比「只是拉格朗日子流形」强得多。注意，由 Liouville–Arnold 定理，拉格朗日环面纤维化与完全可积系统是几乎等价的概念。猜想 2.6.1 主张：任意 Calabi–Yau 流形都具有完全可积系统的结构，而且其守恒量的等位面（level set）是极小子流形。另外，在复 2 维情形，Calabi–Yau 流形是超 Kähler 流形，特殊拉格朗日子流形与全纯子流形通过超 Kähler 旋转联系起来。在这个意义上，SYZ 猜想是椭圆 K3 曲面理论的高维化，可以认为它以某种关于 Calabi–Yau 流形的结构定理为目标。另外，猜想 2.6.3 中出现的 Fourier 变换也称为 SYZ 变换，与族的 Floer 上同调密切相关\footnote{关于族的 Floer 上同调，例如参见 [Fuk02b, Tu15, Aboa, Abob] 等。}。

% [Tutor Note: 上引脚注中文献标签原文作「Aboa, Abob」，疑为 Abo 氏两篇文献的标签写法，照录待核。]

一般认为，与镜像对称相关的诸猜想的强弱关系如下所示：
\[
    \text{位相镜像猜想}\ \Leftarrow\ \text{古典镜像猜想}\ \Leftarrow\ \text{同调镜像猜想}\ \Leftrightarrow\ \text{SYZ 猜想}
\]

由古典镜像对称（Frobenius 流形的同构，或 Hodge 结构的（扩充）形变空间的同构）直接推出位相镜像对称（Hodge 数的对应）。从同调镜像对称推出古典镜像对称，正是 Kontsevich 最初的纲领，虽有种种重要进展，但目前仍未解决\footnote{为从同调镜像对称导出古典镜像对称，需要以下 3 点：1. 从增广三角范畴构成 Frobenius 流形或 [KM94] 意义下的上同调场论（cohomological field theory）；2. 由 1 从 Fukaya 范畴的导出范畴得到 Gromov–Witten 不变量；3. 把 1 应用于相干层的导出范畴，再现 [CdlOGP92] 等人的计算。其中关于 1，有 [Cos07, KKP08, KS09] 等的工作。关于从 Fukaya 范畴计算 Gromov–Witten 不变量，Fukaya 范畴的 Hochschild 上同调与量子上同调的环同构已在 [Kon95b, Sei02] 等中被猜想，并在 [FOOO09, FOOO, Seib, GPS] 等中被研究。另外，3 与斋藤恭司的原始形式理论密切相关 [ST08]。}。若有同调镜像对称，则作为点层的镜像得到拉格朗日环面，并期待它们构成纤维化，这是自然的。作为纤维不取单纯的拉格朗日子流形而取特殊拉格朗日子流形这一猜想，是 Kobayashi–Hitchin 对应（Donaldson–Uhlenbeck–Yau 定理）在辛几何中的类似物，是与拉格朗日子流形稳定性相关的难题。另一方面，从猜想 2.6.3 的主张可以看出，SYZ 猜想包含同调镜像猜想。

镜像对称最初是对 Calabi–Yau 流形表述的，后来被推广到环面 Fano 流形。进而，对于未必是环面的 Fano 流形的镜像对称，也在与环面退化、表示论等相关的各种文献中被讨论 [Giv97, EHX97, Rie08, GS15, Tel]。Fano 流形的镜像对称，看来应在与反典范因子的组（即所谓对数 Calabi–Yau 流形）的文献中来理解 [Aur07, Aur09, GHK15]。另一方面，「对于一般型流形，同调镜像对称也应成立」这一想法经由 [HV00, HIV]，被 Katzarkov 等人所追求 [Kat07, KKOY09, GKR]。受此影响，[Sei11] 证明了亏格 2 曲线的同调镜像对称，并由 [Efi12] 立即推广到一般亏格。另外，对于非紧曲线，亏格 0 的情形由 [AAE$^{+}$13] 证明，并由 [Boc16] 推广到一般亏格。后者的证明本质地使用了二聚体模型。

从 [GP90] 中最初的镜像流形构造也可以看出，Calabi–Yau 流形的镜像对称与 Landau–Ginzburg 模型的镜像对称不可分割地联系在一起。Landau–Ginzburg 模型源自超导的 Ginzburg–Landau 理论\footnote{超导的 Ginzburg–Landau 理论是把二次相变的 Landau 理论应用于超导描述所得的唯象理论，相当于在 Yang–Mills–Higgs 理论中取规范群为 $U(1)$。附带一提，Yang–Mills 理论中作为 BPS 方程的 Hermite–Yang–Mills 方程（Hermitian Yang–Mills equation）在 Yang–Mills–Higgs 理论中的对应物是涡旋方程（vortex equation）。}，但在镜像对称的文献中，它是指由 Kähler 流形及其上的全纯函数的组

% 本 chunk 依原文在 p005 页末截断：c01 自「決まる作用汎関数によって…」接续。
% CHUNK c00 DONE pages p001–p005
